\chapter{SCADA -- Cyber Attacks}
\label{ch:scada_attacks}

\section{Introduzione}
In questo capitolo descriviamo gli scenari d'attacco riprodotti nel testbed, utilizzando come riferimento un dataset Modbus/TCP (baseline).  
Le sperimentazioni sono eseguite esclusivamente in un ambiente di laboratorio isolato e con autorizzazione esplicita; i risultati sono confrontati con studi e dataset pubblici nel dominio SCADA/ICS \cite{scadaAttacks,DatasetModbus}.

\section{Metodo di riproduzione del traffico}
Per generare le condizioni sperimentali abbiamo adottato due modalità principali:
\begin{enumerate}
  \item \textbf{Replay offline:} riproduzione dei pcap originali tramite strumenti come \texttt{tcpreplay} (utile per valutare l'impatto temporale a livello di rete).
  \item \textbf{Rigenerazione programmata:} creazione/alterazione di flussi Modbus con script (ad es. \texttt{Scapy}) per ottenere varianti manipolate, testare payload modificati o pattern temporali specifici.
\end{enumerate}
Con queste modalità abbiamo costruito due condizioni sperimentali:
\begin{description}
  \item[Baseline] traffico originale tratto dal dataset.
  \item[Manipolato] traffico derivato dal baseline e opportunamente alterato per simulare gli attacchi descritti di seguito.
\end{description}

\section{Obiettivi sperimentali}
Gli esperimenti mirano a:
\begin{itemize}
  \item valutare l'impatto funzionale sugli asset ICS (errori Modbus, eccezioni, variazione di valori dei registri);
  \item misurare l'efficacia di rilevamento del sistema \texttt{side-sniff} (NSM/IDS) nel riconoscere anomalie o attacchi;
  \item fornire materiale riproducibile per analisi comparate in laboratorio controllato.
\end{itemize}

\section{Tipologie d'attacco considerate}
Di seguito vengono descritte in modo sintetico e tecnico le tipologie d'attacco simulate nel testbed.

\subsection{Man--In--The--Middle (MITM)}
Un nodo intermedio intercetta e potenzialmente modifica richieste e risposte Modbus/TCP tra master e slave.  
Obiettivi tipici:
\begin{itemize}
  \item lettura non autorizzata dei payload Modbus;
  \item injettare risposte false (modifica dei register values);
  \item introdurre latenza o delay per studiare l'effetto sulla logica applicativa.
\end{itemize}
Nel nostro setup la manipolazione è effettuata offline sui pcap o durante la rigenerazione tramite script per evitare operazioni su sistemi reali.

\subsection{modbusQuery2Flooding}
Variante mirata a inviare coppie di query Modbus ripetute in rapida successione verso lo stesso o più slave.  
Scopo: saturare risorse locali (buffer, CPU, thread) e provocare degrado o malfunzionamenti del servizio Modbus.

\subsection{modbusQueryFlooding}
Flooding generico di query Modbus (letture o scritture) ad alta frequenza; può includere più function code per rendere il pattern meno prevedibile e più difficile da bloccare con semplici rate‑limit.

\subsection{pingFloodDDoS (ICMP)}
Attacco ICMP Echo Request in elevato volume verso l'infrastruttura target. Effetti tipici:
\begin{itemize}
  \item degrado della raggiungibilità della rete di controllo;
  \item aumento dei tempi di risposta ed eventuale perdita di pacchetti.
\end{itemize}

\subsection{tcpSYNFloodDDoS}
Generazione massiva di pacchetti TCP SYN senza completare la handshake, con l'obiettivo di esaurire il backlog di connessioni del target e impedire nuove connessioni legittime (ad esempio verso il servizio Modbus/TCP sulla porta 502).

\section{Organizzazione dei test e note pratiche}
\begin{itemize}
  \item Per ogni attacco si definisce una serie di parametri controllati (tasso di invio, durata, target IP/porta, payload modificati) e si esegue almeno una prova di baseline per confronto.
  \item Tutti i pcap e gli script usati per la rigenerazione/alterazione sono conservati e hashati (SHA256) per garantire riproducibilità e tracciabilità.
  \item Le misure raccolte includono: latenza request--response, errori/exception Modbus, eventuali cambi di stato nei dispositivi simulati e log/alert prodotti da \texttt{side-sniff}.
\end{itemize}

% Citazioni utilizzate nel capitolo
\nocite{scadaAttacks,DatasetModbus}
